% ============================================================================
% Capítulo 2 — Planejamento da solução tecnológica
% ============================================================================
\section{PLANEJAMENTO DA SOLUÇÃO TECNOLÓGICA}

\subsection{Problema a ser resolvido}

O problema central da Cooperativa Vale Verde não é a falta de produção, mas a
falta de coordenação entre produção, demanda e comprovação. A cooperativa sabe
o que seus associados plantam apenas quando alguém responde a uma mensagem; sabe
o que foi entregue apenas quando o motorista retorna; e sabe quanto deve receber
apenas quando o comprador envia a planilha de conferência. Cada uma dessas
lacunas gera retrabalho, perda de margem e, sobretudo, perda de confiança.

A solução proposta, denominada \textbf{AgroConecta}, é uma plataforma corporativa
integrada que registra a cadeia de suprimentos em um único fluxo digital:
cadastro do produtor e sua comunidade, produto, lote de colheita, pedido, rota de
entrega, comprovação com foto e geolocalização, e fatura. O objetivo não é
substituir o cooperativismo, mas dar suporte instrumental para que ele opere com
informação confiável.

\subsection{Objetivos da solução}

O objetivo geral é disponibilizar uma plataforma Web e mobile que integre os
processos de comercialização da agricultura familiar, com rastreabilidade,
segurança e operação offline. Os objetivos específicos estão descritos a seguir.

\begin{itemize}
  \item Centralizar cadastros de cooperativas, produtores, produtos, lotes,
    veículos e motoristas, com controle de acesso por papel.
  \item Registrar pedidos com itens vinculados a lotes, calculando subtotais e
    totais de forma consistente.
  \item Programar rotas que gerem automaticamente as entregas dos pedidos
    selecionados.
  \item Permitir que o motorista atualize o status da entrega e anexe evidência
    fotográfica com coordenadas, mesmo sem conexão.
  \item Emitir faturas a partir dos pedidos, com numeração sequencial e
    idempotência na geração.
  \item Oferecer painéis, relatórios e indicadores para a gestão da cooperativa.
  \item Garantir transparência, acessibilidade e conformidade com a LGPD.
\end{itemize}

\subsection{Público-alvo e personas}

A plataforma atende quatro perfis, cujas necessidades orientaram as decisões de
interface e de permissão. A Tabela~\ref{tab:personas} resume as personas
construídas a partir do levantamento documental e da observação dos processos.

\begin{table}[H]
  \centering
  \caption{Personas e necessidades}
  \label{tab:personas}
  \fontsize{10}{12}\selectfont
  \begin{tabular}{L{2.6cm}L{6.4cm}L{5.4cm}}
    \toprule
    \textbf{Persona} & \textbf{Contexto} & \textbf{Necessidade principal} \\
    \midrule
    Maria, produtora (41) & Mora em comunidade quilombola, cultiva tomate e
      feijão; usa celular Android com internet instável & Registrar a colheita
      no momento da coleta, mesmo sem sinal, e acompanhar a venda \\
    José, motorista (38) & Realiza três rotas por semana em veículo próprio da
      cooperativa & Receber a lista de entregas do dia, atualizar o status e
      comprovar a entrega com foto \\
    Carla, compradora (34) & Nutricionista responsável por compras
      institucionais & Consultar disponibilidade, acompanhar o pedido e obter
      comprovação de entrega \\
    Paulo, gestor (52) & Presidente da cooperativa, acumula funções
      administrativas & Visualizar indicadores, aprovar cadastros e emitir
      faturas com segurança \\
    \bottomrule
  \end{tabular}

  \vspace{2pt}
  {\fontsize{10}{12}\selectfont\itshape Fonte: elaborado pelo autor (2026)}
\end{table}

\subsection{Proposta de valor}

A proposta de valor pode ser sintetizada na seguinte afirmação: \emph{para
cooperativas de agricultura familiar que precisam comprovar origem e entrega,
o AgroConecta é uma plataforma integrada que registra a cadeia produtiva em um
único fluxo digital, com operação offline e evidências verificáveis, reduzindo
perdas e ampliando o acesso a mercados institucionais}. Diferentemente de
sistemas de gestão genéricos, a solução foi desenhada a partir das restrições
reais do campo: conectividade intermitente, aparelhos modestos, usuários com
baixa familiaridade tecnológica e exigência de comprovação documental.

\subsection{Benefícios esperados e diferenciais}

Os benefícios esperados organizam-se em quatro dimensões: \textbf{operacional}
(menos retrabalho na conferência e na programação de rotas), \textbf{comercial}
(maior capacidade de atender compras públicas com documentação consistente),
\textbf{gerencial} (indicadores de vendas, perdas e logística) e \textbf{social}
(inclusão digital de comunidades tradicionais, com interface acessível).

Como diferenciais, destacam-se: operação offline real, com fila de sincronização
persistente; evidência de entrega com foto, coordenadas e vínculo ao pedido;
controle de acesso por papel, com escopo de dados por usuário; e contrato OpenAPI
único, que mantém Web e mobile alinhados à mesma definição de serviço.

\subsection{Viabilidade da solução}

A viabilidade foi analisada sob três aspectos. Do ponto de vista
\textbf{técnico}, a solução utiliza tecnologias maduras e gratuitas no
desenvolvimento (.NET, React Native, SQL Server Developer), com implantação
progressiva. Do ponto de vista \textbf{operacional}, o aplicativo mobile foi
projetado para funcionar em aparelhos Android intermediários e para operar sem
conexão, reduzindo a barreira de adoção. Do ponto de vista \textbf{econômico}, o
custo de infraestrutura em nuvem é compatível com o porte da cooperativa,
conforme estimativa da Tabela~\ref{tab:custos}.

\begin{table}[H]
  \centering
  \caption{Estimativa mensal de infraestrutura em nuvem}
  \label{tab:custos}
  \fontsize{10}{12}\selectfont
  \begin{tabular}{L{4.6cm}L{3.6cm}L{2.6cm}L{3.6cm}}
    \toprule
    \textbf{Componente} & \textbf{Serviço} & \textbf{Estimativa} & \textbf{Observação} \\
    \midrule
    Aplicação Web/API & Azure App Service B1 & R\$ 70--90 & 1 instância, autoescala
      opcional \\
    Banco de dados & Azure SQL Basic & R\$ 25--40 & backup automático \\
    Evidências & Blob Storage & R\$ 5--10 & fotos comprimidas \\
    Segredos & Key Vault & R\$ 3--5 & poucas operações \\
    Observabilidade & Application Insights & R\$ 15--30 & amostragem ajustável \\
    \midrule
    \textbf{Total} & & \textbf{R\$ 120--175} & valores de referência (2026) \\
    \bottomrule
  \end{tabular}

  \vspace{2pt}
  {\fontsize{10}{12}\selectfont\itshape Fonte: elaborado pelo autor (2026), com
  base em preços públicos de referência}
\end{table}

Uma única venda institucional adicional por mês tende a superar o custo mensal
da infraestrutura, o que caracteriza a viabilidade econômica do projeto.

\subsection{Modelo de negócio}

A Tabela~\ref{tab:canvas} apresenta o Lean Canvas da solução, adaptado para o
contexto de uma cooperativa sem fins lucrativos que busca sustentabilidade
operacional.

\begin{table}[H]
  \centering
  \caption{Lean Canvas do AgroConecta}
  \label{tab:canvas}
  \fontsize{9}{11}\selectfont
  \begin{tabular}{L{7.2cm}L{7.2cm}}
    \toprule
    \textbf{Problema} & \textbf{Segmentos de clientes} \\
    Controles manuais, rastreabilidade frágil, perda de produção, baixa
    visibilidade gerencial & Cooperativas de agricultura familiar; produtores
    associados; compradores institucionais \\
    \midrule
    \textbf{Solução} & \textbf{Proposta de valor única} \\
    Plataforma Web + mobile com fluxo produtor--pedido--rota--entrega--fatura e
    operação offline & Rastreabilidade e comprovação de entrega em um único
    fluxo, com inclusão digital \\
    \midrule
    \textbf{Canais} & \textbf{Fontes de receita} \\
    Aplicativo Android, navegador Web, suporte presencial da cooperativa &
    Contribuição dos associados por volume comercializado; economia em perdas \\
    \midrule
    \textbf{Estrutura de custos} & \textbf{Métricas-chave} \\
    Infraestrutura em nuvem, suporte e capacitação & Pedidos por safra, perdas
    por lote, prazo de entrega, produtores ativos \\
    \midrule
    \textbf{Vantagem injusta} & \\
    Conhecimento do território, relação de confiança com as comunidades e
    desenho específico para conectividade intermitente & \\
    \bottomrule
  \end{tabular}

  \vspace{2pt}
  {\fontsize{9}{11}\selectfont\itshape Fonte: elaborado pelo autor (2026),
  adaptado de Osterwalder e Pigneur (2010)}
\end{table}
